% ============================================================
%  APPENDIX A — FedMRI: Educational Federated Learning Platform
%  Place: after \printbibliography in main.tex
%  Figures: copy all PNGs from project root to figures/appendix/
% ============================================================

\appendix
\chapter{FedMRI: An Educational Federated Learning Platform
         for Breast MRI Classification}
\label{app:fedmri}

\section{Overview and Purpose}
\label{app:overview}

FedMRI is an educational web application developed to
demonstrate and validate the federated learning methodology
described in this thesis.
The platform serves a dual purpose: it provides an
interactive showcase of the FedSCRT algorithm
(Chapter~3, Algorithm~\ref{alg:fedscrt}) operating on
real breast DCE-MRI data, and it illustrates the practical
architectural decisions required to deploy federated
learning in a multi-institutional clinical setting
while satisfying the data sovereignty constraints of
Algerian Law~18-07.

The application exposes three distinct user portals
corresponding to three categories of participant in a
federated learning network.
Figure~\ref{fig:fedmri-platform} provides a high-level
overview of the actors and their relationships to the
platform.
The inference engine behind the application is the
trained FedSCRT model: a ConvNeXt-Nano backbone with
Gated Attention MIL head, classifying each uploaded
breast MRI volume as Luminal or Non-Luminal (macro
F1 = 0.629, AUC = 0.687).
Inference is real --- the raw voxel data of each uploaded
\texttt{.mha} or \texttt{.nii} volume is preprocessed
through the identical 64-slice pipeline described in
Section~\ref{sec:preprocessing} and passed through the
PyTorch model.

\begin{figure}[htbp]
  \centering
  \includegraphics[width=0.85\linewidth]%
    {figures/appendix/FedMRI_platform.png}
  \caption{High-level actor diagram of the FedMRI platform.
    Three user roles interact with the central application:
    the Doctor (FL participant, uploads scans and
    contributes to federation training),
    the Patient (FL consumer, benefits from the global
    model without participating in training),
    and the Researcher (FL operator, monitors network
    topology, model metrics, and training logs).
    External integrations include the Flower FL server
    (production), a DICOM viewer for MRI rendering,
    and an AI chatbot for plain-language explanations.}
  \label{fig:fedmri-platform}
\end{figure}


% ─────────────────────────────────────────────────────────
\section{System Architecture}
\label{app:architecture}
% ─────────────────────────────────────────────────────────

The FedMRI platform follows a layered microservices
architecture organised into four tiers
(Figure~\ref{fig:app-architecture}).

\textbf{Frontend layer.}
A Next.js~16 web application (App Router, React~19,
Tailwind CSS, Framer Motion) serves all three portals
under a single codebase with role-based route protection.
A companion Expo mobile application provides the patient
portal on iOS and Android.
All communication to the backend uses JSON over REST and
WebSocket (Socket.io) with JWT tokens validated on every
request.

\textbf{Backend layer.}
A NestJS API (port 3001) implements business logic,
role-based access control, and real-time event emission.
The database layer uses PostgreSQL via Prisma ORM as
the primary store, Redis for session management and
WebSocket pub/sub rate limiting, and MinIO or local
disk for MRI volume storage.
Hospital silo isolation is enforced at the database
query level: every case record carries a
\texttt{scope=HOSPITAL} field that is checked by
middleware on every read, preventing a doctor at one
hospital from accessing another hospital's raw cases.

\textbf{ML service layer.}
A FastAPI service (port 8001) loads the FedSCRT PyTorch
checkpoint and exposes two endpoints:
\texttt{POST /predict} (synchronous, returns subtype,
confidence, and per-slice attention scores) and
\texttt{POST /round/start} (asynchronous, triggers an
FL training round).
The service operates in two modes controlled by an
environment variable: \texttt{INFERENCE\_MODE=real}
(GPU required, live model) and
\texttt{INFERENCE\_MODE=mock} (CPU, deterministic
seeded responses for offline demonstration).

\textbf{FL coordinator.}
A second FastAPI service (port 8002) manages the
federated training round lifecycle.
In development it runs a mock round that replays
pre-computed weight deltas; in production it wraps
\texttt{flwr.server.start\_server()} to orchestrate
real Flower clients on each hospital VM.

\begin{figure}[htbp]
  \centering
  \includegraphics[width=\linewidth]%
    {figures/appendix/app_architechter.png}
  \caption[FedMRI system architecture]{Full system
    architecture of the FedMRI platform.
    The frontend layer (Next.js web + Expo mobile) communicates
    with the NestJS backend via REST and WebSocket.
    The backend dispatches prediction requests synchronously
    to the FastAPI ML service and FL round requests
    asynchronously to the FastAPI FL coordinator.
    The FL coordinator communicates weight deltas only ---
    never raw imaging data --- to the Flower FL server
    in production.
    Storage uses PostgreSQL (primary), Redis
    (sessions + pub/sub), and MinIO/local disk (MRI files).}
  \label{fig:app-architecture}
\end{figure}


% ─────────────────────────────────────────────────────────
\section{Actors and Use Cases}
\label{app:actors}
% ─────────────────────────────────────────────────────────

Figure~\ref{fig:use-cases} presents the complete UML
use case diagram for the three actor roles.
The system boundary separates the Doctor portal,
Patient portal, and Researcher portal.

\begin{figure}[htbp]
  \centering
  \includegraphics[width=0.72\linewidth]%
    {figures/appendix/fedmri_use_case_diagram 1.png}
  \caption{UML use case diagram for the FedMRI platform.
    The Doctor actor interacts with the system as an FL
    participant: uploading scans triggers both immediate
    inference and an asynchronous FL training round.
    The Patient actor interacts as an FL consumer: scanning
    and receiving predictions in plain language, with no
    exposure to FL terminology or model internals.
    The Researcher actor has read-only access to aggregate
    model metrics, network topology, and training logs,
    but never to raw patient data.}
  \label{fig:use-cases}
\end{figure}

\subsection{Doctor Portal (FL Participant)}
\label{app:doctor}

The doctor is the primary FL participant.
Each clinical interaction follows the same flow:
(1)~upload a breast MRI scan;
(2)~receive an immediate binary subtype prediction
    (Luminal / Non-Luminal), a confidence score,
    and a per-slice attention heatmap identifying
    the slices most informative to the model decision;
(3)~validate or dispute the prediction (active learning
    feedback loop).

A critical architectural constraint is that the
doctor's result is returned \emph{before} the FL
round fires.
The FL round executes asynchronously in the background,
as shown in the sequence diagram (Figure~\ref{fig:doctor-sequence}).
This design ensures that clinical response time is
bounded by the inference latency (typically 2--5 seconds
for a single MRI volume on a GPU) and not by the
federation training time (~30 seconds in mock mode,
longer in production).

The doctor portal provides seven screens: Dashboard,
Scan Analysis, Medical History, AI Assistant,
Model Performance, Documentation, and Settings.
Hospital silo isolation is enforced throughout:
a doctor at Hospital~A can view only Hospital~A's
cases.

\begin{figure}[htbp]
  \centering
  \includegraphics[width=0.82\linewidth]%
    {figures/appendix/FedMRI___Sequence_Doctor_Upload__FL_Round_AutoTrigger.png}
  \caption[Doctor upload and FL round sequence]{Sequence
    diagram for the Doctor scan upload
    and FL round auto-trigger flow.
    The doctor submits a scan (JWT authenticated).
    The NestJS backend forwards the image path to the
    ML service synchronously and returns the prediction
    result and attention map to the doctor (solid arrows).
    An FL round is then triggered asynchronously: the FL
    coordinator sequentially notifies each of the three
    hospital clients (Hospital~A, B, C), collects weight
    deltas, and emits a WebSocket \texttt{fl:round:complete}
    event when the round finishes.
    At no point are raw MRI images transmitted across
    the hospital boundary (\texttt{rawDataTransmitted=0}
    is recorded in the privacy audit log).}
  \label{fig:doctor-sequence}
\end{figure}

\subsection{Patient Portal (FL Consumer)}
\label{app:patient}

The patient is an independent consumer of the global
federated model with no hospital affiliation and no
participation in the training loop.
Patient-facing language deliberately avoids FL
terminology: the platform communicates
``AI trained across 3 hospitals'' rather than
``federated learning'' or ``gradient aggregation''.

The patient uploads their own MRI file and receives
the binary classification result in plain clinical
language, accompanied by an advisory statement:
Luminal subtypes are associated with hormone-receptor
positivity and may indicate eligibility for endocrine
therapy; all results require clinical correlation and
specialist review.
The AI assistant (backed by Claude via the Anthropic
API, streamed via Server-Sent Events) provides
plain-language explanations of the result on request.
The patient may download a PDF summary of their
analysis session.

\begin{figure}[htbp]
  \centering
  \includegraphics[width=0.72\linewidth]%
    {figures/appendix/FedMRI___Sequence_Patient_MRI_Upload.png}
  \caption{Sequence diagram for the patient MRI upload flow.
    The patient's scan is processed by the same FastAPI ML
    service as the doctor's scan, but the resulting case
    record carries \texttt{scope=PATIENT}, isolating it
    from the hospital silo entirely.
    No FL round is triggered by a patient upload: patients
    are consumers of the trained model, not contributors
    to its training.}
  \label{fig:patient-sequence}
\end{figure}

\subsection{Researcher Portal (FL Operator)}
\label{app:researcher}

The researcher has a global read-only view of the
federation: model performance metrics across rounds,
network topology showing the three hospital nodes and
their training status, aggregated dataset statistics,
and training logs.
The researcher never accesses raw patient images or
individual case records.
This access control is enforced at the middleware level:
all researcher queries return aggregate or model-level
data, never row-level patient data.

\begin{figure}[htbp]
  \centering
  \includegraphics[width=0.72\linewidth]%
    {figures/appendix/FedMRI___Sequence_Researcher_Dataset_Access.png}
  \caption{Sequence diagram for the researcher dataset
    access flow.
    The researcher requests dataset statistics and model
    metrics via the NestJS API.
    The response contains aggregated counts, F1 trajectories,
    and weight delta norms --- never raw case images or
    patient identifiers.
    The privacy guarantee holds by construction: the
    researcher role has no database permissions on the
    \texttt{Case} table beyond count queries.}
  \label{fig:researcher-sequence}
\end{figure}


% ─────────────────────────────────────────────────────────
\section{Database Schema}
\label{app:database}
% ─────────────────────────────────────────────────────────

Figure~\ref{fig:db-schema} presents the Prisma entity
relationship diagram for the FedMRI database.
Nine entities model the complete application domain.

\textbf{Hospital} and \textbf{User} model the
institutional and clinical actors.
Each user belongs to exactly one hospital, and each
hospital has a unique Flower client identifier for
the federated training protocol.

\textbf{Case} is the central entity.
Its \texttt{scope} field enforces the silo boundary:
a \texttt{scope=HOSPITAL} case is visible only to the
submitting hospital; a \texttt{scope=PATIENT} case is
visible only to the submitting patient.
The \texttt{storedLocally=true} invariant ensures that
image files never leave the institution that uploaded them.

\textbf{FLRound} and \textbf{FLContribution} record the
federated training history.
\texttt{FLContribution.rawDataTransmitted} is a
database-level invariant that is always zero, enforcing
the FL privacy claim in a verifiable, auditable way.

\textbf{PrivacyAuditLog} records every FL round event
with the hospital identifier, event type, bytes
transmitted, and a \texttt{rawDataTransmitted} field
that is always zero --- providing the audit trail
required by Law~18-07 article~22 for health data processing.

\textbf{Feedback} captures active learning signals:
when a doctor disputes a prediction, the corrected
subtype label is stored and the \texttt{alTriggered}
flag indicates whether an additional FL round was
automatically initiated.

\begin{figure}[htbp]
  \centering
  \includegraphics[width=\linewidth]%
    {figures/appendix/FedMRI___database_schema.png}
  \caption[FedMRI database entity-relationship diagram]{%
    Entity-relationship diagram for the FedMRI
    PostgreSQL schema (Prisma ORM).
    Nine entities cover hospital and user management,
    case storage with silo isolation, FL round
    lifecycle tracking, active learning feedback,
    AI chat messages, model metrics, and privacy
    audit logging.
    The \texttt{rawDataTransmitted=0} field in both
    \texttt{FLContribution} and \texttt{PrivacyAuditLog}
    is a database-level invariant that enforces the
    federated privacy guarantee.}
  \label{fig:db-schema}
\end{figure}


% ─────────────────────────────────────────────────────────
\section{Deployment Architecture}
\label{app:deployment}
% ─────────────────────────────────────────────────────────

The platform supports two deployment configurations
controlled entirely by environment variables, requiring
zero application code changes to switch between them.

\subsection{Development: Single Machine, Docker Compose}

The development configuration
(Figure~\ref{fig:dev-deployment}) runs all services
on a single machine using Docker Compose.
The ML service operates with
\texttt{INFERENCE\_MODE=mock} (no GPU required) and
the FL coordinator with \texttt{FL\_MODE=mock}
(replays pre-computed weight deltas rather than
running live Flower clients).
This configuration allows the full application
including UI, authentication, FL round animation,
and prediction results to be demonstrated on
any laptop without GPU hardware.

\begin{figure}[htbp]
  \centering
  \includegraphics[width=0.70\linewidth]%
    {figures/appendix/Mock_Dev__Single_Machine__Docker_Compose.png}
  \caption{Development deployment on a single machine
    via Docker Compose.
    Five containers: Next.js web frontend (:3000),
    NestJS API (:3001), FastAPI ML service (:8001,
    mock mode), FastAPI FL coordinator (:8002, mock
    mode), PostgreSQL (:5432), Redis (:6379), and
    local disk storage.
    No GPU is required; all FL rounds replay seeded
    deterministic responses.}
  \label{fig:dev-deployment}
\end{figure}

\subsection{Production: Multiple VMs, Real Flower}

The production configuration
(Figure~\ref{fig:prod-deployment}) distributes services
across three virtual machines.
VM~1 hosts the web application and NestJS API behind
an Nginx reverse proxy.
VM~2 hosts the FastAPI ML service with
\texttt{INFERENCE\_MODE=real}, requiring a GPU for
live inference.
VM~3 hosts the Flower FL coordinator with
\texttt{FL\_MODE=flower}, running
\texttt{flwr.server.start\_server()} and managing
three Flower clients deployed on separate hospital VMs
(VM~4, VM~5, VM~6).

Critically, each hospital VM (4--6) holds only its
local patient data.
The Flower protocol transmits model weight deltas only
across the hospital firewall boundary.
Raw MRI volumes, patient identifiers, and label
information never traverse the hospital network
boundary, satisfying the data sovereignty requirements
of Law~18-07 and the federated privacy model of
Chapter~2.

\begin{figure}[htbp]
  \centering
  \includegraphics[width=0.90\linewidth]%
    {figures/appendix/Production__Multiple_VMs.png}
  \caption{Production deployment across six virtual machines.
    VM~1 (app server): Nginx + Next.js + NestJS API.
    VM~2 (ML server, GPU required): FastAPI inference
    with real FedSCRT checkpoint.
    VM~3 (FL coordinator): Flower server
    (\texttt{flwr.server.start\_server()}).
    VMs~4--6 (hospital clients): one per simulated
    hospital, each holding local data only.
    The Flower protocol carries weight deltas only across
    the hospital firewalls; the label \emph{``weight deltas
    only cross boundary''} on each hospital connection
    reflects the FL privacy invariant.}
  \label{fig:prod-deployment}
\end{figure}


% ─────────────────────────────────────────────────────────
\section{Privacy Guarantees by Design}
\label{app:privacy}
% ─────────────────────────────────────────────────────────

The FedMRI platform implements the federated privacy
model at four independent layers, such that the
privacy guarantee holds by construction and cannot
be violated by application logic errors.

\textbf{Layer 1 — Data localisation.}
MRI files are stored on the local disk of the hospital
server that uploaded them.
The \texttt{Case.storedLocally=true} database invariant
is enforced by the NestJS service layer on every case
creation.
The ML service receives only a file path, never raw
pixel data, over the internal network.

\textbf{Layer 2 — Silo query isolation.}
Every database query on the \texttt{Case} table
includes a \texttt{WHERE hospitalId = :current\_hospital}
filter injected by NestJS middleware.
No application code can read another hospital's cases
without deliberately bypassing this middleware.

\textbf{Layer 3 — FL protocol constraint.}
The Flower protocol transmits only serialised PyTorch
model state dictionaries between the coordinator and
the hospital clients.
No patient data, no image tensors, and no label
information are present in the Flower message payloads.

\textbf{Layer 4 — Audit log invariant.}
Every FL round completion writes a
\texttt{PrivacyAuditLog} row with
\texttt{rawDataTransmitted=0}.
This field is set by the FL coordinator service and
is never modified by application logic.
The audit log provides the verifiable evidence trail
required by Law~18-07 article~22 for health data
processing by automated systems.

Table~\ref{tab:privacy-layers} summarises the four
layers and their enforcement mechanisms.

\begin{table}[htbp]
  \centering
  \caption{Privacy guarantee layers in the FedMRI
    platform. Each layer independently enforces the
    data sovereignty constraint; all four must be
    bypassed simultaneously to violate the guarantee.}
  \label{tab:privacy-layers}
  \renewcommand{\arraystretch}{1.3}
  \begin{tabular}{llll}
    \toprule
    \textbf{Layer} & \textbf{Mechanism}
      & \textbf{Enforcement} & \textbf{Failure mode} \\
    \midrule
    Data localisation
      & \texttt{storedLocally=true}
      & Service layer invariant
      & File system error \\
    Silo isolation
      & Hospital ID middleware filter
      & NestJS guard
      & Auth bypass \\
    FL protocol
      & Weight deltas only
      & Flower serialisation
      & Protocol modification \\
    Audit log
      & \texttt{rawDataTransmitted=0}
      & DB write on round close
      & DB corruption \\
    \bottomrule
  \end{tabular}
\end{table}


% ─────────────────────────────────────────────────────────
\section{Technology Stack Summary}
\label{app:stack}
% ─────────────────────────────────────────────────────────

Table~\ref{tab:tech-stack} summarises the complete
technology stack of the FedMRI platform.

\begin{table}[htbp]
  \centering
  \caption{FedMRI platform technology stack.}
  \label{tab:tech-stack}
  \renewcommand{\arraystretch}{1.3}
  \begin{tabular}{lll}
    \toprule
    \textbf{Component} & \textbf{Technology}
      & \textbf{Role} \\
    \midrule
    Web frontend    & Next.js 16, React 19, Tailwind CSS
      & Doctor + Researcher portals \\
    Mobile frontend & Expo React Native
      & Patient portal (iOS/Android) \\
    API backend     & NestJS, Prisma ORM
      & Business logic, auth, events \\
    Primary DB      & PostgreSQL
      & Persistent data store \\
    Cache/pub-sub   & Redis
      & Sessions, WebSocket events \\
    File storage    & MinIO / local disk
      & MRI volume storage \\
    ML service      & FastAPI, PyTorch
      & Inference + attention maps \\
    FL coordinator  & FastAPI, Flower (flwr)
      & Training round management \\
    AI assistant    & Anthropic Claude API
      & Plain-language explanations \\
    Real-time       & Socket.io, SSE
      & FL round progress events \\
    \bottomrule
  \end{tabular}
\end{table}


% ─────────────────────────────────────────────────────────
\section{Application Screenshots}
\label{app:screenshots}
% ─────────────────────────────────────────────────────────

The following screenshots illustrate the FedMRI web application
as rendered in a desktop browser.
All screens were captured from the development deployment
running with \texttt{INFERENCE\_MODE=mock} and
\texttt{FL\_MODE=mock}.

% ── Authentication ────────────────────────────────────────

\begin{figure}[htbp]
  \centering
  \begin{subfigure}[t]{0.48\linewidth}
    \includegraphics[width=\linewidth]%
      {figures/appendix/app-scrennshots/FedMRI_Login.png}
    \caption{Login screen.}
    \label{fig:ss-login}
  \end{subfigure}
  \hfill
  \begin{subfigure}[t]{0.48\linewidth}
    \includegraphics[width=\linewidth]%
      {figures/appendix/app-scrennshots/FedMRI_Signup.png}
    \caption{Sign-up screen.}
    \label{fig:ss-signup}
  \end{subfigure}
  \caption{Authentication screens. Users register with a role
    (Doctor, Patient, or Researcher) and a hospital affiliation.
    JWT tokens are issued on successful login and validated by
    NestJS middleware on every subsequent request.}
  \label{fig:ss-auth}
\end{figure}

% ── Doctor portal ─────────────────────────────────────────

\begin{figure}[htbp]
  \centering
  \includegraphics[width=0.90\linewidth]%
    {figures/appendix/app-scrennshots/Doctor_Portal_Dashboard.png}
  \caption{Doctor portal --- Dashboard.
    Displays the hospital's aggregated case statistics,
    recent scan activity, and the current federated model
    version.
    The sidebar exposes the seven doctor screens:
    Dashboard, Scan Analysis, Medical History,
    AI Assistant, Model Performance,
    Documentation, and Settings.}
  \label{fig:ss-doctor-dashboard}
\end{figure}

\begin{figure}[htbp]
  \centering
  \includegraphics[width=0.90\linewidth]%
    {figures/appendix/app-scrennshots/New_MRI_Scan_Upload_Analysis_Flow.png}
  \caption{Doctor portal --- Scan Upload and Analysis.
    The doctor uploads a \texttt{.mha} or \texttt{.nii}
    MRI volume; the FastAPI ML service returns a binary
    subtype prediction (Luminal / Non-Luminal),
    a confidence score, and a per-slice attention heatmap
    within 2--5 seconds.
    An FL training round is triggered asynchronously
    in the background.}
  \label{fig:ss-scan-upload}
\end{figure}

\begin{figure}[htbp]
  \centering
  \includegraphics[width=0.90\linewidth]%
    {figures/appendix/app-scrennshots/Medical_History_Archive.png}
  \caption{Doctor portal --- Medical History Archive.
    Lists all cases submitted by the current hospital,
    with subtype prediction, confidence, submission date,
    and validation status.
    Hospital silo isolation ensures that only cases with
    \texttt{hospitalId} matching the authenticated user
    are returned.}
  \label{fig:ss-medical-history}
\end{figure}

\begin{figure}[htbp]
  \centering
  \includegraphics[width=0.90\linewidth]%
    {figures/appendix/app-scrennshots/Doctor_Portal_Documentation.png}
  \caption{Doctor portal --- Documentation.
    Provides inline reference documentation for the
    FedSCRT model architecture, the federated training
    protocol, and the interpretation of attention heatmaps.
    Content is rendered from Markdown and does not
    require an external network request.}
  \label{fig:ss-doctor-docs}
\end{figure}

% ── Patient portal ────────────────────────────────────────

\begin{figure}[htbp]
  \centering
  \includegraphics[width=0.90\linewidth]%
    {figures/appendix/app-scrennshots/Ask_AI_Patient_Chat.png}
  \caption{Patient portal --- AI Assistant chat.
    The patient submits a natural-language question about
    their scan result.
    Responses are streamed from the Anthropic Claude API
    via Server-Sent Events and rendered incrementally in
    the chat window.
    FL terminology is deliberately absent from
    patient-facing language.}
  \label{fig:ss-patient-chat}
\end{figure}

% ── Researcher portal ─────────────────────────────────────

\begin{figure}[htbp]
  \centering
  \begin{subfigure}[t]{0.48\linewidth}
    \includegraphics[width=\linewidth]%
      {figures/appendix/app-scrennshots/Researcher_Portal_Dataset_Management.png}
    \caption{Dataset Management.}
    \label{fig:ss-dataset-mgmt}
  \end{subfigure}
  \hfill
  \begin{subfigure}[t]{0.48\linewidth}
    \includegraphics[width=\linewidth]%
      {figures/appendix/app-scrennshots/dataset_ready.png}
    \caption{Dataset ready confirmation.}
    \label{fig:ss-dataset-ready}
  \end{subfigure}
  \caption{Researcher portal --- Dataset Management.
    Left: the researcher inspects per-hospital case counts,
    class distributions, and partition statistics.
    Right: confirmation banner indicating that the dataset
    partition is balanced across the three hospital silos
    and that the FL coordinator is ready to initiate a
    new training round.}
  \label{fig:ss-researcher-dataset}
\end{figure}

\begin{figure}[htbp]
  \centering
  \begin{subfigure}[t]{0.48\linewidth}
    \includegraphics[width=\linewidth]%
      {figures/appendix/app-scrennshots/Researcher_Portal_Model_Performance.png}
    \caption{Model Performance.}
    \label{fig:ss-model-perf}
  \end{subfigure}
  \hfill
  \begin{subfigure}[t]{0.48\linewidth}
    \includegraphics[width=\linewidth]%
      {figures/appendix/app-scrennshots/FL_insights.png}
    \caption{FL Insights.}
    \label{fig:ss-fl-insights}
  \end{subfigure}
  \caption{Researcher portal --- Model Performance and FL Insights.
    Left: macro F1, AUC, and per-class precision/recall
    curves across federated rounds.
    Right: FL Insights panel showing per-hospital weight
    delta norms, round duration, and participation status
    for the most recent training round.}
  \label{fig:ss-researcher-perf}
\end{figure}

\begin{figure}[htbp]
  \centering
  \includegraphics[width=0.90\linewidth]%
    {figures/appendix/app-scrennshots/Researcher_Portal_System_Logs.png}
  \caption{Researcher portal --- System Logs.
    Chronological event log for all FL round lifecycle
    events: round start, client notifications, weight
    delta reception, aggregation, and round completion.
    Each row records the hospital identifier, event type,
    bytes transmitted, and the \texttt{rawDataTransmitted}
    field (always zero).}
  \label{fig:ss-system-logs}
\end{figure}

\begin{figure}[htbp]
  \centering
  \includegraphics[width=0.90\linewidth]%
    {figures/appendix/app-scrennshots/Html_Body.png}
  \caption{Application shell --- HTML body structure.
    Browser developer tools view of the rendered DOM root,
    illustrating the Next.js App Router hydration structure
    and the Tailwind CSS utility class pattern applied
    throughout the platform.}
  \label{fig:ss-html-body}
\end{figure}
